\section{Conclusion}\label{sec:short_conclusion}

This paper quantifies three consequences of volatility estimation for WTI Average Price Options: the point-price correction from posterior integration, the posterior dispersion of conditional prices, and the pricing effect of changing volatility information. Holding the posterior and pricing model fixed isolates the integration correction. Its local dependence on posterior variance and volatility curvature explains why it grows in weak-information, high-curvature states and why partial fixing can alter its magnitude.

In the observed October-2026 panel, integration changes the point price little even though posterior model-price intervals remain appreciably wider than that correction. Across later-date valuations, earlier option-implied volatility produces much lower settlement-pricing errors than the historical inputs evaluated here. Recent historical estimates and horizon-matched GARCH do not reproduce that performance; same-contract IV persistence and external vanilla-WTI information do. The empirical finding concerns volatility information under a fixed valuation model, rather than a structural volatility-risk-premium decomposition.

The resulting diagnostic extends the interpretation of small integration effects: closeness of PI and PM establishes little movement in the point value, not little uncertainty around it. For the evaluated WTI samples, the larger point-pricing effect comes from the volatility information supplied to the model. This conclusion remains conditional on the common-factor specification, available settlements, and limited number of market dates.
